\documentclass[10pt,a4paper]{article}
\usepackage[utf8]{inputenc}
\usepackage[T1]{fontenc}
\usepackage{lmodern}
\usepackage[margin=1.8cm,top=2.0cm,bottom=2.0cm]{geometry}
\usepackage{amsmath,amssymb,amsfonts}
\usepackage{graphicx}
\usepackage{xcolor}
\usepackage{tcolorbox}
\tcbuselibrary{skins,breakable}
\usepackage{enumitem}
\usepackage{booktabs}
\usepackage{adjustbox}
\usepackage{microtype}
\usepackage{hyperref}

% Color definitions
\definecolor{brandblue}{HTML}{2C5AA0}
\definecolor{brandgreen}{HTML}{2E7D52}
\definecolor{brandamber}{HTML}{C8881E}
\definecolor{brandred}{HTML}{B23A48}
\definecolor{brandpurple}{HTML}{6A4C93}
\definecolor{brandink}{HTML}{21355E}
\definecolor{bgfill}{HTML}{EAF0F7}

% Custom tcolorbox environments
\newtcolorbox{intuition}[1][Intuition]{
  enhanced, breakable, colback=bgfill, colframe=brandblue,
  fonttitle=\bfseries\color{white}, title={#1},
  boxrule=0.8pt, arc=3pt, left=10pt, right=10pt, top=8pt, bottom=8pt
}

\newtcolorbox{worked}[1][Worked Example]{
  enhanced, breakable, colback=white, colframe=brandgreen,
  fonttitle=\bfseries\color{white}, title={#1},
  boxrule=0.8pt, arc=3pt, left=10pt, right=10pt, top=8pt, bottom=8pt
}

\newtcolorbox{everyday}[1][Everyday Picture]{
  enhanced, breakable, colback=bgfill!50, colframe=brandamber,
  fonttitle=\bfseries\color{white}, title={#1},
  boxrule=0.8pt, arc=3pt, left=10pt, right=10pt, top=8pt, bottom=8pt
}

\newtcolorbox{watchout}[1][Watch Out]{
  enhanced, breakable, colback=white, colframe=brandred,
  fonttitle=\bfseries\color{white}, title={#1},
  boxrule=0.8pt, arc=3pt, left=10pt, right=10pt, top=8pt, bottom=8pt
}

\newtcolorbox{keytake}[1][Key Takeaways]{
  enhanced, breakable, colback=bgfill, colframe=brandpurple,
  fonttitle=\bfseries\color{white}, title={#1},
  boxrule=0.8pt, arc=3pt, left=10pt, right=10pt, top=8pt, bottom=8pt
}

\newenvironment{steps}{
  \begin{enumerate}[label=\textbf{Step \arabic*:}, leftmargin=*]
}{
  \end{enumerate}
}

\newcommand{\housefig}[2]{
  \begin{center}
    \includegraphics[width=0.88\textwidth]{#1}\\
    \small\textcolor{brandink}{\textbf{Figure:} #2}
  \end{center}
}

\setlength{\parindent}{0pt}
\setlength{\parskip}{6pt}

\begin{document}

% Banner Header
\begin{tcolorbox}[enhanced, colback=brandink, colframe=brandink, arc=0pt, left=12pt, right=12pt, top=12pt, bottom=12pt]
  \color{white}
  \large\textbf{ISM Companion Reader} \hfill \textbf{Session 9}\\[4pt]
  \huge\textbf{Testing of Hypothesis Foundations}\\[4pt]
  \small Null & alternative hypotheses, Type I & II errors, p-values, Z-tests for means and proportions.
\end{tcolorbox}

\vspace{10pt}

\section{Principles of Hypothesis Testing}

Hypothesis testing is a structured statistical framework used to validate claims about population parameters using sample data.

\begin{intuition}[The Courtroom Analogy]
In a criminal trial, the defendant is assumed innocent until proven guilty beyond a reasonable doubt.
\begin{itemize}
  \item \textbf{Null Hypothesis ($H_0$):} Defendant is innocent (status quo / no effect).
  \item \textbf{Alternative Hypothesis ($H_1$):} Defendant is guilty (the claim we seek to prove).
  \item \textbf{Burden of Proof:} We only reject $H_0$ if sample evidence is overwhelmingly unlikely under $H_0$.
\end{itemize}
\end{intuition}

\section{Type I and Type II Errors}

Because sample data carries randomness, statistical decisions can lead to two types of errors:
\begin{itemize}
  \item \textbf{Type I Error ($\alpha$):} Rejecting $H_0$ when $H_0$ is actually true (False Alarm).
  \item \textbf{Type II Error ($\beta$):} Failing to reject $H_0$ when $H_1$ is true (Missed Detection).
  \item \textbf{Power of the Test ($1 - \beta$):} Probability of correctly rejecting a false $H_0$.
\end{itemize}

\housefig{figures/fig_errors.png}{Trade-offs between Type I Error $\alpha$, Type II Error $\beta$, and Test Power $1 - \beta$.}

\section{One-Tailed vs Two-Tailed Z-Tests}

Depending on the alternative hypothesis, rejection regions sit in one tail or split across both tails.

\begin{align*}
\text{Two-Tailed Test:} \quad & H_0: \mu = \mu_0 \quad \text{vs} \quad H_1: \mu \neq \mu_0 \\
\text{Right-Tailed Test:} \quad & H_0: \mu \le \mu_0 \quad \text{vs} \quad H_1: \mu > \mu_0 \\
\text{Left-Tailed Test:} \quad & H_0: \mu \ge \mu_0 \quad \text{vs} \quad H_1: \mu < \mu_0
\end{align*}

\housefig{figures/fig_two_tailed.png}{Two-tailed Z-test with critical bounds $Z = \pm 1.96$ at significance level $\alpha = 0.05$.}

\housefig{figures/fig_right_tailed.png}{Right-tailed Z-test with critical bound $Z = +1.645$ at significance level $\alpha = 0.05$.}

\begin{worked}[Slide Example: Testing Difference of Two Population Means]
\textbf{Problem Statement:} Two machine processes produce steel rods. A sample of $n_1 = 50$ rods from Machine 1 yields mean breaking strength $\bar{X}_1 = 38$ kg with $\sigma_1 = 4$ kg. A sample of $n_2 = 60$ rods from Machine 2 yields $\bar{X}_2 = 35$ kg with $\sigma_2 = 5$ kg. Test if there is a significant difference between the true population means at $\alpha = 0.05$.

\begin{steps}
\item \textbf{State Hypotheses:}
\[
H_0: \mu_1 - \mu_2 = 0 \quad \text{vs} \quad H_1: \mu_1 - \mu_2 \neq 0
\]
\item \textbf{Determine Critical Value:} For a two-tailed Z-test at $\alpha = 0.05$, critical values are $Z_{\text{crit}} = \pm 1.96$.
\item \textbf{Calculate Standard Error of Difference:}
\[
\text{SE}_{\bar{X}_1 - \bar{X}_2} = \sqrt{\frac{\sigma_1^2}{n_1} + \frac{\sigma_2^2}{n_2}} = \sqrt{\frac{4^2}{50} + \frac{5^2}{60}} = \sqrt{\frac{16}{50} + \frac{25}{60}} = \sqrt{0.320 + 0.417} = \sqrt{0.737} \approx 0.858
\]
\item \textbf{Compute Z Test Statistic:}
\[
Z_{\text{cal}} = \frac{(\bar{X}_1 - \bar{X}_2) - 0}{\text{SE}} = \frac{38 - 35}{0.858} = \frac{3}{0.858} \approx 3.50
\]
\item \textbf{Make Decision:} Since $Z_{\text{cal}} = 3.50 > 1.96$, we reject $H_0$. There is strong evidence of a significant difference between the two machines.
\end{steps}
\end{worked}

\begin{worked}[Slide Example: Test of Significance for Single Proportion]
\textbf{Problem Statement:} A phone survey claims that 60\% ($p_0 = 0.60$) of voters favor candidate A. In a sample of $n = 400$ voters, $X = 220$ favor candidate A ($\bar{p} = 220/400 = 0.55$). Test if candidate support is significantly lower than claimed at $\alpha = 0.05$.

\begin{steps}
\item \textbf{State Hypotheses:}
\[
H_0: p \ge 0.60 \quad \text{vs} \quad H_1: p < 0.60 \quad (\text{Left-Tailed Test})
\]
\item \textbf{Determine Critical Value:} For left-tailed test at $\alpha = 0.05$, $Z_{\text{crit}} = -1.645$.
\item \textbf{Calculate Standard Error:}
\[
\text{SE}_{\bar{p}} = \sqrt{\frac{p_0(1-p_0)}{n}} = \sqrt{\frac{0.60 \times 0.40}{400}} = \sqrt{\frac{0.24}{400}} = \sqrt{0.0006} \approx 0.0245
\]
\item \textbf{Compute Z Test Statistic:}
\[
Z_{\text{cal}} = \frac{\bar{p} - p_0}{\text{SE}_{\bar{p}}} = \frac{0.55 - 0.60}{0.0245} = \frac{-0.05}{0.0245} \approx -2.04
\]
\item \textbf{Make Decision:} Since $Z_{\text{cal}} = -2.04 < -1.645$, we reject $H_0$. Support is significantly lower than 60\%.
\end{steps}
\end{worked}

\begin{watchout}[p-Value Decision Rule]
Always compare $p$-value directly to significance level $\alpha$:
\begin{itemize}
  \item If $p\text{-value} \le \alpha$, \textbf{Reject $H_0$} ("If $p$ is low, $H_0$ must go").
  \item If $p\text{-value} > \alpha$, \textbf{Fail to reject $H_0$}.
\end{itemize}
\end{watchout}

\begin{keytake}[Summary Checklist]
\begin{itemize}
  \item \textbf{$H_0$ vs $H_1$:} $H_0$ assumes equality/no effect; $H_1$ holds the claim to prove.
  \item \textbf{Type I Error $\alpha$:} False alarm probability set by the analyst.
  \item \textbf{Decision Rule:} Reject $H_0$ when $|Z_{\text{cal}}| > Z_{\text{crit}}$ or $p\text{-value} \le \alpha$.
\end{itemize}
\end{keytake}

\end{document}
